Nanopore sekvenovanie je technológia používaná na sekvenovanie DNA a RNA molekúl \cite{Zheng2023Nanopore}. Táto metóda umožňuje čítanie genetického materiálu tým, že molekuly prechádzajú cez nanometrové póry (nanopóry), čo vedie k zmene elektrického prúdu, ktorý je následne analyzovaný na určenie sekvencie nukleotidov \cite{Zheng2023Nanopore}.
Je to revolučná technológia v oblasti genomiky, ktorá umožňuje rýchle a presné sekvenovanie s minimálnymi nákladmi a zariadeniami prenosných veľkostí \cite{Zheng2023Nanopore}.
Aplikácie nanopórového sekvenovania zahŕňajú diagnostiku chorôb, výskum genetických porúch a monitorovanie environmentálnych vzoriek \cite{Zheng2023Nanopore}.
